\chapter{LITERATURE REVIEW}

\section{Visual Geolocation in Mountainous Terrain}

Visual geolocation estimates the location where a photograph was taken using information contained in the image. In urban environments, buildings, roads, and other man-made structures provide many distinctive landmarks that can be matched with reference images. Mountain environments are much more challenging because natural landscapes contain fewer unique visual features. Their appearance also changes with weather, lighting, snow cover, and seasonal vegetation, making conventional image matching less reliable.

Researchers found that one feature remains relatively stable despite these changes: the mountain skyline. Although snow and vegetation may change throughout the year, the boundary between the mountains and the sky is determined by the terrain itself. This makes the skyline a reliable geographic feature for estimating camera position. Early studies demonstrated that mountain skylines generated from \Glspl{dem} could be aligned with photographs to recover the camera location and orientation \cite{naval1997,gupta2008camera}. These works established the foundation for later skyline-based localization methods.

A major step towards large-scale visual geolocation was presented by \cite{baatz2012}. Instead of matching individual mountain peaks, they represented panoramic skylines using contour descriptors extracted from \Gls{dem}-generated terrain models. Since a photograph usually contains only part of the surrounding horizon, the system searched for reference panoramas that contained the observed skyline rather than requiring a complete match. Their experiments over approximately 40,000 km$^2$ of Switzerland demonstrated that skyline shape alone contains sufficient information for large-scale localization.

Building on this work, \cite{saurer2016image} observed that similar skyline shapes could sometimes appear at different locations, especially when only a small portion of the skyline was visible. To reduce false matches, they improved both skyline extraction and candidate verification. After retrieving the most likely candidate locations, they compared the complete skyline with the rendered terrain model before selecting the final match. This additional verification step significantly improved localization accuracy while maintaining efficient retrieval.

These studies established skyline matching as a practical approach for mountain geolocation. They also demonstrated the value of \Glspl{dem}, which make it possible to generate reference skylines without requiring photographs from every possible location. As the field matured, researchers shifted their attention from demonstrating the feasibility of skyline matching to improving its robustness under more challenging imaging conditions and different types of terrain.

\section{\Glspl{dem}}

A \Gls{dem} is a digital representation of the Earth's terrain. In skyline-based geolocation, a \Gls{dem} is used to generate synthetic views of the landscape from different virtual camera positions. These rendered views provide the reference skylines that are later compared with photographs. As a result, the quality of the \Gls{dem} directly affects the accuracy of the generated skyline.

Several global \Glspl{dem} are available for mountainous terrain, including NASADEM, ALOS AW3D30, and Copernicus GLO-30. Each has strengths depending on the application. NASADEM generally provides better ground elevation in densely forested regions because its radar measurements penetrate vegetation more effectively. AW3D30 performs well in some rugged mountain landscapes but may introduce artificial terrain artefacts. Copernicus GLO-30 provides a more detailed representation of ridges and valleys, making it well suited for applications that depend on terrain morphology rather than vegetation structure \cite{Li2022}.

The Khumbu region is dominated by steep mountain terrain where the overall shape of ridges and skylines is more important than modelling the ground beneath dense vegetation. Since this research generates synthetic skylines from terrain data, preserving ridge geometry is the primary requirement. Previous evaluations have shown that Copernicus GLO-30 provides detailed and consistent representations of mountainous terrain, making it suitable for skyline generation and visual geolocation \cite{Li2022}.

For these reasons, Copernicus GLO-30 was selected to generate the synthetic mountain views and reference skylines used throughout this research.

% -------------------------------------------------------------------------
% Alternative version (if using the \Glspl{hma} 8 m \Gls{dem})
%
% The Khumbu region lies within \Glspl{hma}, where high-resolution
% \Glspl{dem} are available for glacierized terrain. The \Glspl{hma}
% 8 m \Gls{dem} provides substantially finer spatial resolution than global
% 30 m \Glspl{dem} and was developed specifically for analysing glacier geometry
% and elevation change. Its higher resolution preserves fine ridge
% structure and terrain details that may be smoothed in coarser datasets
% \cite{Guo2023}.
%
% Since skyline generation depends primarily on the accurate representation
% of mountain ridges rather than vegetation or land cover, the \Gls{hma} 8 m \Gls{dem}
% offers an attractive alternative for this study. Its detailed terrain
% representation is particularly suitable for the rugged landscapes of the
% Khumbu region, where sharp ridges and steep slopes dominate the skyline.
%
% For these reasons, the \Glspl{hma} 8 m \Gls{dem} can be used to generate
% the synthetic mountain views and reference skylines employed in this
% research.
% -------------------------------------------------------------------------

\section{Improving Skyline Matching}

Although skyline matching proved to be effective, researchers soon identified several practical limitations. In real-world photographs, the skyline is not always clearly visible. Clouds, haze, vegetation, and foreground objects may partially hide the mountain boundary. In addition, different locations can sometimes produce similar skyline shapes, making it difficult to distinguish between nearby viewpoints. These challenges motivated the development of more robust skyline representations and matching techniques.

One important limitation was identified by \cite{chen2015camera}. Earlier methods relied primarily on the outer skyline, assuming that it could always be extracted reliably from the input image. However, the authors showed that weather conditions, image quality, and atmospheric effects often make the skyline incomplete or unreliable. To overcome this problem, they incorporated mountain ridge lines in addition to the skyline. Ridge lines were extracted from depth discontinuities in \Gls{dem}-generated terrain models and provided additional geometric information when the skyline alone was insufficient. The method also estimated the camera roll angle during localization, showing that even small rotations could noticeably reduce matching accuracy if ignored. By combining skyline and ridge information, the system produced more reliable localization under challenging viewing conditions.

While \cite{chen2015camera} improved the robustness of skyline matching, their method was developed primarily for large mountain ranges where the skyline changes gradually with camera movement. A different challenge arises in hilly landscapes, where the skyline is much closer to the observer. In these environments, even small changes in camera position can significantly alter the shape of the skyline.

To address this problem, \cite{pan2020camera} proposed a skyline representation designed specifically for nearby hills. Instead of using only the outer mountain boundary, they introduced \emph{lapel points}, which are formed where internal ridge lines intersect the skyline. These points act as stable landmarks that make neighbouring skylines easier to distinguish. The authors also compared skylines at multiple image scales, allowing the system to compensate for small differences in viewing distance that would otherwise lead to incorrect matches. Their results demonstrated that incorporating additional terrain structure significantly improved localization accuracy in hilly environments.

As skyline databases became larger, the computational cost of searching every reference skyline also increased. Rather than improving the skyline representation itself, \cite{pan2021fast} focused on making the retrieval process more efficient. Similar skylines were first grouped into clusters, reducing the number of candidate matches that needed to be examined. The authors also showed that the geometric relationship between matched skyline features could be used to estimate the camera orientation directly from the image, reducing the dependence on external orientation sensors. These improvements reduced localization time while maintaining comparable accuracy.

Although these methods advanced skyline matching considerably, they were developed for landscapes that differ from the Himalayan environment. The European studies focused on the Alps, while the work of \cite{pan2021fast} considered densely hilly terrain where nearby viewpoints produce rapidly changing skylines. In the Himalayas, mountain peaks are typically much farther from the observer, causing the skyline to remain relatively stable over neighbouring viewpoints. Consequently, techniques designed specifically for nearby hills, such as dense lapel-point representations and fine-scale viewpoint sampling, become less important. Instead, robust skyline extraction and efficient matching are more critical for reliable localization in high mountain regions.

A notable benchmark for visual geolocation in mountainous terrain is the GeoPose3K dataset \cite{brejcha2019geopose3k}, which contains mountain images with associated ground truth locations. This dataset has been used by multiple research groups to evaluate and compare localization methods under real-world conditions, making it useful for situating new approaches within the broader field.

\section{Skyline Extraction}

The success of any skyline-based localization system depends on how accurately the skyline can be extracted from an input image. Even if a reference database and matching algorithm are highly accurate, errors introduced during skyline extraction can propagate through the rest of the localization pipeline. In practice, extracting the skyline is often challenging because outdoor images contain clouds, haze, trees, buildings, and varying illumination that obscure the true boundary between the mountains and the sky.

One of the earliest studies to address this problem was presented by \cite{woo2007robust}. The authors observed that conventional edge detectors faced a trade-off between detection and localization. A detector tuned to identify only strong edges often missed parts of the skyline, while one tuned to capture fine details also detected many unwanted edges produced by clouds, shadows, and vegetation. To overcome this limitation, they generated edge maps at two different scales. A coarse-scale edge map was used to identify reliable skyline points, while a finer-scale edge map traced the complete skyline between these points. This two-stage approach produced more accurate skyline extraction under difficult outdoor conditions.

Although edge-based methods improved skyline detection, they relied heavily on carefully selected image-processing parameters. Their performance often decreased when weather conditions or lighting changed significantly. Rather than manually designing filters to detect the skyline, later research explored learning-based approaches that could adapt automatically to different image characteristics.

A resource-efficient learning-based method was proposed by \cite{ahmad2021resource}. Instead of using a deep neural network, the authors trained a bank of lightweight linear filters to distinguish skyline edges from other image edges. During inference, the local image structure determined which filter should be applied at each pixel, allowing the method to suppress unwanted edges while preserving the mountain boundary. A dynamic programming algorithm was then used to recover a continuous skyline from the filtered responses. This approach achieved accuracy comparable to more complex learning-based methods while requiring considerably less computational power, making it suitable for resource-constrained platforms such as mobile devices, unmanned air vehicles, and planetary rovers.

Despite these improvements, both edge-based and shallow learning methods still treated skyline extraction primarily as an edge detection problem. Recent advances in computer vision instead formulate skyline extraction as a semantic segmentation task. Rather than identifying individual edges, semantic segmentation classifies every pixel in the image as either sky or terrain. The skyline is then obtained directly from the boundary between the two regions. Since the prediction is based on information from the entire image instead of local gradients alone, semantic segmentation is generally more robust to clouds, shadows, vegetation, and varying illumination.

This shift towards semantic segmentation has enabled more reliable skyline extraction in challenging mountain environments. Encoder-decoder architectures such as U-Net combine high-level semantic information with fine image details, allowing accurate pixel-level segmentation while preserving object boundaries. These properties make semantic segmentation particularly suitable for mountain geolocation, where precise skyline boundaries are essential for reliable matching with \Gls{dem}-generated reference skylines.

\section{Synthetic Data for Mountain Geolocation}

As visual geolocation methods became increasingly dependent on machine learning, the availability of labelled training data emerged as a major challenge. Unlike urban environments, remote mountain regions contain relatively few publicly available photographs, and manually annotating mountain skylines is both time-consuming and labour-intensive. This lack of training data limits the performance of learning-based skyline extraction and localization methods.

To address this problem, \cite{tomesek2022crosslocate} proposed CrossLocate, a cross-modal visual geolocation framework that combines real photographs with synthetic terrain rendered from \Glspl{dem}. Instead of comparing only real images, the authors rendered depth maps and terrain views from many virtual camera positions. A deep neural network was then trained to learn a shared feature representation between rendered terrain and real photographs. This approach allowed the system to localize images even in regions where very few real training photographs were available.

An important finding of their work was that synthetic data is not simply a replacement for real images. The authors showed that there is a noticeable gap between synthetic and real image domains, meaning that models trained only on one type of data do not generalize well to the other. By jointly learning from both domains, CrossLocate achieved significantly better localization performance over large mountainous regions than methods relying on a single data source.

More recently, \cite{liu2023cmlocate} extended this idea through a complete cross-modal localization framework designed for natural environments without \Gls{gnss} information. Their system combined semantic segmentation, \Gls{dem}-generated reference data, and efficient database retrieval to localize images directly from terrain features. Rather than relying solely on skyline contours, the framework learned richer semantic representations while maintaining practical retrieval speeds for large reference databases. Their results demonstrated that combining semantic understanding with terrain information can improve both localization accuracy and retrieval efficiency.

These studies demonstrate that synthetic terrain generated from \Glspl{dem} provides an effective solution to the limited availability of labelled mountain imagery. Synthetic data makes it possible to generate training samples and reference views for locations where few or no photographs exist. This is particularly valuable in remote mountain regions such as the Himalayas, where collecting large, well-annotated image datasets is difficult. Consequently, synthetic data has become an important component of modern mountain geolocation systems.

\section{Comparison of Existing Methods}

Table~\ref{tab:literature_comparison} summarizes the main approaches discussed in this chapter. Early studies focused on demonstrating that mountain skylines could be used for localization. Later research improved skyline matching, developed more robust skyline extraction methods, and introduced synthetic terrain to overcome the limited availability of mountain image datasets. Together, these studies established skyline-based localization as a practical approach for natural environments.

\begin{table}[H]
\centering
\caption{Comparison of representative mountain geolocation methods.}
\label{tab:literature_comparison}
\small
\begin{tabular}{p{2.8cm}p{2.8cm}p{2.5cm}p{3.4cm}}
\toprule
\textbf{Study} & \textbf{Main Contribution} & \textbf{Terrain} & \textbf{Limitation} \\
\midrule
Baatz et al.~\cite{baatz2012}
&
Contour-based skyline matching
&
Alps
&
Requires reliable skyline extraction
\\

Saurer et al.~\cite{saurer2016image}
&
Improved skyline verification
&
Alps
&
Designed mainly for Alpine terrain
\\

Chen et al.~\cite{chen2015camera}
&
Skyline and ridge matching
&
Mountain regions
&
Higher computational complexity
\\

Pan et al.~\cite{pan2020camera}
&
Lapel points for nearby hills
&
Hilly terrain
&
Assumes rapidly changing skylines
\\

Pan et al.~\cite{pan2021fast}
&
Efficient skyline retrieval
&
Hilly terrain
&
Sensitive to pseudo skylines
\\

Yang et al.~\cite{woo2007robust}
&
Robust edge-based skyline extraction
&
Mountain regions
&
Sensitive to image conditions
\\

Ahmad et al.~\cite{ahmad2021resource}
&
Lightweight skyline extraction
&
Mountain regions
&
Still relies on edge responses
\\

Tomesek et al.~\cite{tomesek2022crosslocate}
&
Cross-modal learning using synthetic terrain
&
Large-scale mountains
&
Requires deep learning and large datasets
\\

Liu et al.~\cite{liu2023cmlocate}
&
Cross-modal semantic localization
&
Natural environments
&
Computationally intensive
\\
\bottomrule
\end{tabular}
\end{table}

\section{Research Gap}

Previous research has shown that mountain skylines provide reliable information for visual geolocation. Researchers have improved skyline matching, developed more robust skyline extraction methods, and demonstrated the usefulness of synthetic terrain generated from \Glspl{dem}. Together, these advances have established skyline-based localization as a practical alternative when conventional image matching is unreliable.

Despite this progress, most existing systems have been developed for environments that differ from the Himalayas. Many studies focus on the European Alps, where large image datasets are available and mountain viewpoints are relatively well documented. Other work targets nearby hilly terrain, where small changes in camera position produce significant changes in skyline shape. These assumptions do not necessarily hold in the Khumbu region, where distant mountain ranges create more stable skyline profiles over neighbouring viewpoints.

Another limitation is that many recent approaches rely on computationally expensive deep learning models or large image databases. While these methods achieve high localization accuracy, they are less suitable for lightweight deployment on mobile devices operating in remote mountain environments. In addition, satellite navigation may be unreliable in deep valleys where surrounding terrain blocks \Gls{gnss} signals, making image-based localization an attractive alternative.

These observations highlight the need for a localization system designed specifically for the Himalayan environment. Such a system should operate without internet connectivity, generate reference data directly from terrain models, and perform efficient skyline matching while remaining practical for deployment on resource-constrained devices. This research addresses these requirements by combining semantic skyline extraction, synthetic terrain generation, and hierarchical skyline matching for visual geolocation in the Khumbu region.
